\chapter{线性代数与多重线性代数。}

线性代数既是数学中最古老的分支之一，又是最新近的分支之一。一方面，在数学的起源处就可找到一些通过单次乘法或除法即可解决的问题，也就是说，通过计算函数 \(f(x) = ax\) 的一个值，或通过解方程 \(ax = b\) 来解决：这些正是线性代数的典型问题，而不“线性地”思考，就不可能处理它们，甚至不可能正确地陈述它们。

另一方面，不仅这些问题，而且几乎一切涉及一次方程的内容，长期以来都被贬入初等教学，直到域、环、拓扑向量空间等概念的现代发展，才解放了线性代数的那些本质概念（例如对偶），并赋予它们重要性；正是在那时，人们认识到了几乎整个现代代数在本质上是线性的——这种“线性化”甚至是其显著特征之一——线性代数也重新获得了它应有的地位。因此，从我们所采取的观点出发来叙述它的历史，将是一项既重要又困难的任务；我们在这里只能满足于相当简略的提示。

由前文所述可知，线性代数无疑是为了满足从事实际计算者的需要而诞生的；因此我们看到，三率法（{[}311{]}，第 I 卷，第~150-155~页）与试位法以或多或少明确的形式，在一切实用算术教科书中扮演着重要角色，从埃及人的莱因德纸草书，经阿耶波多、阿拉伯人、比萨的莱昂纳多以及中世纪与文艺复兴时期不可胜数的“算书”，直到今日我们小学里仍在使用的教科书；但它们或许从来不过是供从业者之用、从更高级的科学理论中摘出的一部分。

至于真正的数学家，他们关于线性代数的研究的性质，乃是其科学的一般结构的函数。古希腊数学，如欧几里得《原本》中所阐述的那样，发展了两个具有线性特征的抽象理论：一方面是量的理论（{[}107{]}，第 V 卷；参见第~152~页），另一方面是整数的理论（{[}107{]}，第 VII 卷）。在巴比伦人那里，我们发现了一些与我们初等代数接近得多的方法；他们知道如何求解一次方程组，而且方式十分优雅（{[}232{]}，第~181-183~页）。然而，在很长一段时期内，线性代数的进步主要是与代数运算的进步联系在一起的，而正是从这种为本注记所不涉及的角度去考察它才是合理的；事实上，要把一个线性方程组化归为 \(ax = b\) 型的方程，当只涉及一个未知数时，只需知道那些（丢番图已经实质上陈述过的）把各项从一个等号边移到另一边并合并同类项的规则就够了；而当涉及多个未知数时，只需再知道如何逐个消去它们，直到只剩下一个即可。因此，直到 18 世纪的那些代数论著都认为，一旦列出了这些规则，在一次方程方面就已经做完了一切；至于未知数个数与方程个数相同的方程组（他们不考虑别的情形），当其各方程的左端不是线性无关的型时，他们总是满足于顺带指出这意味着问题提得不恰当。在 19 世纪的论著中，甚至在某些更晚近的著作中，这一观点也只是靠记号的进步——它使写出 \(n\) 个未知数的 \(n\) 个方程的方程组成为可能——以及行列式的引入——它使在“一般情形”下陈述显式求解的公式成为可能——才有所改变；这一进步，倘若莱布尼茨（{[}198 a{]}，第 I 卷，第~239~页）曾发表过他在这方面的想法，本应归功于他，而它实际上主要归功于 18 世纪及 19 世纪初的数学家们。

但是，我们必须首先考察几股思想潮流，它们对线性代数——按我们所理解的意义——的发展所作的贡献，远比对线性方程组的研究为大。费马受阿波罗尼奥斯著作的启发（{[}109{]}，第 I 卷，第~91-110~页；法文译文，第 III 卷，第~85-101~页），甚至在笛卡儿（{[}85 a{]}，第 VI 卷）之前就构思了解析几何的原理，想到了借助次数对平面曲线进行分类（这一分类到 17 世纪末已为所有数学家越来越熟悉，可以认为已确定地被掌握），并陈述了基本原理：在平面上，一次方程表示一条直线，二次方程表示一条圆锥曲线：他由此立即推出了一些关于几何轨迹的“非常漂亮的”结论。与此同时，他预告了（{[}109{]}，第 I 卷，第~184-188~页；法文译文，第 III 卷，第~159-163~页）对问题的这样一种分类：已解决的问题；可化归为含两个未知数的方程的问题；可化归为含三个未知数的方程的问题；等等；他还补充道：第一类在于确定一个点，第二类在于确定一条直线或一个平面轨迹，其后各类在于确定一个曲面；等等。（“\ldots 这样一个问题所寻求的不仅是一个点或一条线，而是一个与问题相称的整个曲面；正是由此产生了那些面轨迹，其后各类依此类推”，出处同上，第~186~页；在那里已经可以找到 n 维几何的胚胎。）这段陈述了代数中和几何代数中的维数原理的文字，显示出了代数与几何的一种融合，它与现代观念完全合拍，然而如前所见，它花了两个多世纪才深入人心。

至少这些思想不久便导致了解析几何的兴盛，它在 18 世纪随着克莱罗、欧拉、克拉默、拉格朗日以及许多其他人而获得其全部规模。平面和空间中坐标变换公式的线性特征——费马不可能不曾注意到它——例如被欧拉（{[}108 a{]}，（1），第 IX 卷，第 II-III 章及附录第 IV 章）加以突出，他在此基础之上建立了平面曲线以及曲面按次数的分类（正是因为这些公式的线性性，次数才是不变的）；也是他（出处同上，第 XVIII 章）引入了“仿射性”一词，用以指称可以经由变换 \(x' = ax, y' = by\) 互相推演出的那些曲线之间的关系（但在这一定义中，他没有察觉到任何几何不变量，因为该定义仍然依附于坐标轴的一个特定选取）。稍后，我们看到拉格朗日（{[}191{]}，第 III 卷，第~661-692~页）以一整部论文——它将长期理所当然地享有盛名——专门研究三维解析几何中典型的线性和多重线性问题。大约在同一时期，联系到求过若干给定点的平面曲线这一线性问题，行列式的概念成形了，最初以某种多少是经验的方式，见于克拉默 {[}72{]} 和贝祖 {[}21{]}；这一概念随后被多位作者加以发展，其本身及其本质性质最终由柯西（{[}56 a{]}，（2），第 I 卷，第~91-169~页）和雅可比（{[}171{]}，第 III 卷，第~355-392~页）明确地提炼了出来。

另一方面，当数学家们对一次方程还多少有些轻视之际，微分方程的求解却相反构成一个重大问题；在这些方程之中，人们自然很早就挑出了线性方程——不论其系数是否为常数——对它们的研究有助于提升线性性及与之相伴诸物的价值。达朗贝尔 {[}75 c{]}、拉格朗日（{[}191{]}，第 I 卷，第~471-668~页）和欧拉（{[}108 a{]}，（1），第 XII 卷）的情形正是如此；但其中只有第一位认为值得说出：非齐次方程的通解是一个特解与相应的齐次方程的通解之和；此外，当这些作者断言 n 阶线性齐次方程的通解是 n 个特解的线性组合时，他们并没有补充说这些特解必须是线性无关的，也不曾费力把后一概念弄明确；在这一点上，正如在许多其他问题上一样，看来只有柯西在综合理工学校的教学才澄清了问题（{[}56 b{]}，第~573-574~页）。但是拉格朗日（出处同上）已经同样引入了（诚然仅仅是为了计算，而且没有给它命名）伴随方程 \(L^{*}(y) = 0\)，即线性微分方程 \(L(y) = 0\) 的伴随，这是对偶性的一个典型例子，其依据是关系式

\[
\int z L (y) d x = \int L ^ {*} (z) y d x
\]

它在 y 与 z 于积分区间两端为零时成立；更确切地说，在高斯明确给出 3 个变量的线性代换的转置之前 30 年，我们在这里无疑看到了“函数算子”的第一个例子：\(L^{*}\) 是借助一个双线性函数（这里是积分 \(\int yzdx\)）给出的算子 L 的转置或“伴随”。

与此同时，同样是在拉格朗日那里（{[}191{]}，第 III 卷，第~695-795~页），线性代换——起初是 2 个或 3 个变量的——开始向算术进军。显然，函数 \(F(x, y)\) 当 \(x\) 与 \(y\) 取遍全部整数值时所取的值组成的集合，在对 \(x, y\) 施行任意一个具整系数且行列式为 1 的线性代换时并不改变；拉格朗日正是以这一基本观察为基础，建立了数用型表示的理论以及型的化简理论；而高斯则以一步——我们如今已难以完全估量其全部大胆之处——由此提炼出等价的概念和型的类的概念（参见第~50~页）；在这一点上，他认识到需要若干关于线性代换的初等原理，并特别引入了转置或伴随的概念（{[}124 a{]}，第 I 卷，第~304~页）。从那时起，2 个、3 个以及后来的 \(n\) 个变量的二次型的算术研究和代数研究，与它们紧密相连的双线性型的研究，以及更近时期这些概念向无穷多个变量的推广，直到我们这个时代，一直是线性代数最富有成果的进步源泉之一（参见第~136~页以下）。

但是，或许更为决定性的一步是，高斯在这部《算术研究》中创立了（参见第~51~页）有限阿贝尔群的理论，它以四种不同的方式出现：整数关于模 m（m 为整数）的加法群的方式，与 m 互素的数关于模 m 的乘法群的方式，二元二次型的类群的方式，以及最后，m 次单位根的乘法群的方式；而且，如我们已经指出的，高斯显然是把所有这些群当作阿贝尔群、或者更确切地说当作 Z 上的模来处理的，研究它们的结构、它们的同构关系等等。在“复整数”\(a + bi\) 的模中，我们后来看到他研究一个 Z 上的无限模，他无疑没有忽略它与椭圆函数的周期模（由他在复域中发现）的同构；无论如何，这一思想在雅可比那里已经清楚地出现，例如在他关于三个周期的函数不可能性的著名证明中，以及在他关于阿贝尔积分反演问题的见解中（{[}171{]}，第 II 卷，第~25-50~页），并很快以克罗内克的诸定理告终（{[}186 a{]}，第 III\(_{1}\) 卷，第~49-109~页）。

在这里，在我们一直试图追寻其脉络、有时还有其迂回的这几股潮流之中，又有另一股长期潜行于地下的潮流汇入。正如别处将更详细地表明的那样（见第~130~页），前一世纪所理解意义下的“纯粹”几何学，也就是说，本质上是不用坐标的平面和空间的射影几何，是 17 世纪由德萨格 {[}84{]} 创立的，他的思想曾被一位费马赏识其真正价值、被一位帕斯卡付诸使用，随后便归于废弃，被解析几何的辉煌进步所遮蔽；到 18 世纪末，它由蒙日，随后由彭赛列及其追随者布里昂雄、沙勒，重新带回荣耀的地位，有时完全地、有意地与解析方法相分离，有时（特别是在德国）又与后者紧密地混合在一起。然而，射影变换，无论从哪种观点（综合的或解析的）加以考虑，当然无非是射影坐标或“重心”坐标上的线性代换；圆锥曲线（在 17 世纪）、以及后来的二次曲面——它们的射影理论长期是这一学派的主要研究对象——无非就是二次型，而我们上文已经揭示过二次型与线性代数的密切联系。与这些概念相汇合的还有配极的概念：极点与极线的理论同样由德萨格创立，它在蒙日及其后继者手中、不久又以对偶原理之名，成为变换几何定理的一种有力工具；即使人们不敢断言它与伴随微分方程的联系只是很晚才被注意到（平谢莱在世纪末指出了这种联系），有一点毕竟没有被错过——沙勒对此作证（{[}60 b{]}，第~55~页）——那就是注意到它与球面三角形对偶（互反球面三角形）概念的亲缘关系，后者早在 16 世纪就由韦达（{[}319{]}，第~418~页）和斯涅耳引入球面三角学。但是射影几何中的对偶性只是向量空间对偶性的一个方面，这时要考虑到从仿射空间过渡到射影空间（它是仿射空间对“数乘”关系的商空间）所强加的那些修正。

19 世纪比我们历史上的任何时期都更富有一流的数学家；而要在寥寥几页之中——即使只限于最突出的方面——描述这些思想运动在他们手中的融合所产生的一切，是很困难的。一方面是纯粹综合的方法，它像一个普罗克鲁斯提斯之床，其正统信徒们在其中自我折磨，另一方面是与任意强加于空间的坐标系相联系的解析方法，人们很快就感到需要一种几何演算，它被莱布尼茨梦想过而未创立，被卡诺 {[}51{]} 不完善地勾勒过：首先出现的是向量的加法，它隐含于高斯对虚数的几何表示及他由此对初等几何所作的应用之中（参见第~162~页），由贝拉维蒂斯以“等偏差法”之名加以发展，并在格拉斯曼、默比乌斯、哈密顿那里取得其确定的形式；与此同时，默比乌斯以“重心演算”之名给出了它的一个适于射影几何之需的版本（{[}223{]}，第 I 卷）。

与此同时，在同样一些人中间，实现了从平面和“通常”空间到 n 维空间的过渡，这一过渡是如此自然（一旦走上这条道路），以至于我们已经看到它被费马预告过；这一过渡甚至是不可避免的，因为对于 2 个或 3 个变量可用几何术语自身来解释的代数现象，对任意多个变量都毫无改变地成立；因此，若在使用几何语言时强加一个限于 2 维或 3 维的限制，对现代数学家来说将是一个桎梏，其碍事程度不亚于那个一直阻止希腊人把数的概念扩张到不可公度之量的比的桎梏。也正因为如此，与 n 维空间有关的语言和观念几乎同时在各方面出现，在高斯那里是隐约的，在下一代数学家那里则是清晰的；而他们在使用这些观念时或大或小的大胆程度，与其说取决于他们的数学倾向，不如说取决于他们的哲学见解、甚至仅仅实际的考虑。无论如何，凯莱和格拉斯曼在大约 1846 年就以极大的自如驾驭这些概念（凯莱在与格拉斯曼相异之处说道（{[}58{]}，第 I 卷，第~321~页），“无需诉诸任何形而上学的概念”）；在凯莱那里，人们始终非常接近于解析解释和坐标，而在格拉斯曼那里，从一开始——随着 n 维空间中向量的加法——领先的是几何方面，并最终引出我们稍后将要谈到的那些发展。

然而，来自高斯的推动正以两种不同的方式把数学家引向代数即超复系统的研究。一方面，人们不可避免地要尝试以引入“虚数单位”\(i = \sqrt{-1}\) 之外的其他方式来扩张实数的范围，从而也许能开辟出比复数域更广阔、更富成果的领域。高斯使自己确信（{[}124 a{]}，第 II 卷，第~178~页）这种扩张是不可能的，至少当人们想保留复数的主要性质、也就是说用现代的语言讲、使它成为一个交换域的那些性质时是如此；而他的同时代人，或由于他的影响，或独立地，似乎也分享了这一信念，它直到很久以后才由魏尔斯特拉斯（{[}329 a{]}，第 II 卷，第~311-332~页）以一个精确的定理加以证明。但是，一旦把复数的乘法解释为平面中的旋转，人们若想把这一想法推广到空间，就会被引导去考虑（由于空间中的旋转构成一个非交换群）非交换的乘法；这正是引导哈密顿 \(^{1}\) 发现四元数 {[}145 a{]} 的思想之一，四元数是除环的第一个例子。这一例子的独特性（正如弗罗贝尼乌斯后来所证明的，它是在实数域上可以构造出来的唯一例子）在某种程度上限制了它的影响范围，尽管有、或者甚至恰恰由于形成了一个狂热的“四元数主义者”学派：一种奇怪的现象，它后来在格拉斯曼的著作周围重演，随后又有那些普及者从哈密顿和格拉斯曼那里提炼出所谓“向量演算”。稍后格雷夫斯和凯莱放弃结合律而构造出“凯莱数”，也并未开辟出一条有意义的矿脉。但是，在西尔维斯特引入了矩阵并（未给它命名）明确定义了秩（{[}304{]}，第 I 卷，第~145-151~页）之后，又是凯莱（{[}58{]}，第 II 卷，第~475-496~页）创立了矩阵的演算，而且他不曾忽略（这是一个本质的事实，后来常被忽视）指出：矩阵无非是线性代换的一种缩写记号，正如事实上高斯用 \((a, b, c)\) 表示型 \(aX^2 + 2bXY + cY^2\) 一样。这里的情形除了是西尔维斯特和凯莱关于行列式及与之相关的一切的大量论著的一个方面——对我们来说无疑是最有趣的一个方面——之外，别无其他，那些论著布满了巧妙的恒等式和令人印象深刻的计算。

格拉斯曼发现的（除其他东西外）还有实数域上的一个代数，即与他的名字联系在一起的外代数。他的著作（{[}134{]}，第 I₁ 卷）甚至先于哈密顿的著作，是在一种几乎完全彻底的精神孤独中创造出来的，长期鲜为人知，这无疑是由于它的独创性，也由于它开头笼罩着的哲学迷雾——例如这迷雾一开始就使默比乌斯望而却步。格拉斯曼怀着与哈密顿类似但视界更为宽广的关切（而且他很快就意识到，这些关切正是莱布尼茨本人的关切），建造起一座庞大的代数—几何大厦，它奠基于 n 维向量空间的一种几何的或者说“内蕴”的观念（差不多已经公理化）之上；在他最终得到的结果之中，最初等的一些包括：向量线性无关的定义、维数的定义以及基本关系式 dim V + dim W = dim (V + W) + dim (V ∩ W)（出处同上，第~209~页；参见{[}134{]}，第 I₂ 卷，第~21~页）。但首先是多向量的外乘法、继之内乘法，为他提供了工具，使他得以从容地处理线性代数本身的问题，然后是涉及欧氏结构即向量正交性的那些问题（在后者中他找到了他所缺少的对偶性的等价物）。

高斯在超复系统研究中开辟的另一条道路，是从复整数 \(a + bi\) 出发的那一条；从这些整数出发，人们很自然地继续走向整数环 Z 上或有理数域 Q 上的更一般的代数或超复系统，首先走向高斯已经预见到的、由单位根生成的那些系统，然后走向代数数域和代数整数模：这些正是库默尔著作的主要对象，而对前者的研究则由狄利克雷、埃尔米特、克罗内克、戴德金着手进行。在这里，与实数域上的代数的情形相反，无须放弃交换域的任何一条特征性质，而整个 19 世纪的工作也正是限于这些系统。但那些线性的性质，例如寻求域中整数的基（它对于判别式的一般定义是不可或缺的），在许多方面都扮演着本质的角色；而且，无论如何在戴德金那里，那些方法注定要变成典型“超复的”；此外，戴德金本人虽然并未以一般的方式处理代数的问题，却意识到他工作的这一方面，意识到例如把它与魏尔斯特拉斯关于实数域上超复系统的结果联系起来的东西（{[}79{]}，特别是第 II 卷，第~1-19~页）。与此同时，代数数域中单位乘法群结构的确定——由狄利克雷在其著名的论文（{[}92{]}，第 I 卷，第~619-644~页）以及几乎同时由埃尔米特（{[}159{]}，第 I 卷，第~159-160~页）完成——极适合于澄清关于 Z 上的模、它们的生成系以及（当它们具有时）它们的基的观念。随后是理想的概念：戴德金在代数数域中把它定义为（域的整数环上的）一个模（{[}79{]}，第 III 卷，第~251~页），而克罗内克则在多项式环中（以“模系”之名）引入一个等价的概念（{[}186 a{]}，第 II 卷，第~334-342~页），这给出了 Z 以外更一般的环上的模的第一批例子；而且，在这些相同的作者那里、然后在希尔伯特那里，带算子群的概念以及在特殊情形下从这样一个群出发总能构造出一个适当定义的环上的模的可能性，被一点一点地释放出来。

与此同时，对二次型和双线性型的算术—代数研究及其“化简”（或者换一种说法，矩阵及其“不变量”的化简），导致发现了关于线性方程组求解的一般原理，这些原理曾因缺少秩的概念而未被雅可比注意到。\(^{2}\) 整系数线性方程组的整数解问题被提出并解决，起初是由埃尔米特在一个特殊情形中，然后由 H. J. 史密斯以其全部一般性解决（{[}287{]}，第 I 卷，第~367-409~页）；后者的结果直到 1878 年才由弗罗贝尼乌斯重新发现，那是在克罗内克所制定、魏尔斯特拉斯也参与其中的一个庞大的研究纲领的框架之内；顺便说一句，正是克罗内克在研究过程中给出了实（或复）系数线性方程组诸定理的最终形式，这些定理也由《爱丽丝漫游奇境记》的著名作者在一部不起眼的教材中以他所特有的缜密细心加以阐明；至于克罗内克，他不屑于发表自己的结果，而把它留给同事和弟子们；甚至连“秩”这个词也只是由弗罗贝尼乌斯引入的。同样是在柏林大学的教学中，克罗内克 {[}186 b{]} 和魏尔斯特拉斯给出了行列式的“公理化”定义（即 n 维空间中 n 个向量的一个多线性反对称函数，规范化为在单位矩阵上取值 1），这一定义等价于可以从格拉斯曼的计算中推导出来的那个定义；也还是在他的讲义中，克罗内克在仍未感到需要给它命名、且尚非内蕴形式的情况下，引入了空间的张量积和矩阵的“克罗内克”积（由施于各因子的给定线性代换在张量积上诱导出的线性代换）。

这些研究也不应与不变量理论相分离，后者由凯莱、埃尔米特和西尔维斯特创立（埃尔米特后来在书信中称之为“不变量三位一体”），而从现代的观点看，它首先是一种线性群的表示理论。作为射影几何对偶性的代数等价物，出现了同步变量列与逆步变量列的区分，也就是说，一个空间中的向量与对偶空间中的向量的区分；在首先注意到 2 个或 3 个变量的低次型、然后是任意次数的型之后，人们很快便考察起多组“同步”或“逆步”变量的双线性型、继而多线性型来，这等价于张量的引入；当里奇和列维—奇维塔在不变量理论的启发下于 1900 年把“张量分析”引入微分几何 {[}257{]} 时，这一点便变成了自觉的并得到普及——张量分析后来因“相对论者”物理学家们的使用而风靡一时。也还是不变量理论、微分几何以及偏微分方程理论（特别是所谓普法夫问题及其推广）的逐步相互渗透，逐渐引导几何学家们首先考虑微分的双线性反对称型，特别是一次型的“双线性协变量”（1870 年由李普希茨引入，随后由弗罗贝尼乌斯研究），最终导致 E. 嘉当（{[}52 a{]}，第 II₁ 卷，第~303-396~页）和庞加莱（{[}251 b{]}，第 III 卷，第 XXII 章）创立外微分形式的演算。庞加莱是在构成其积分不变量的脉络中，把它们作为出现在多重积分中的表达式引入的，而嘉当——无疑是在他关于代数的研究的引导之下——则以一种不那么形式的方式引入它们，但也不曾放过指出其演算的代数部分与格拉斯曼的外乘法完全相同（他采用的名字即由此而来），从而最终把后者的著作放回到它真正的位置上。把外微分形式翻译成张量分析的记号，此外还立即表明了它们与反对称张量的联系，而一旦回到纯代数的观点，这就会使人看出：它们之于交错多线性型，正如协变张量之于任意多线性型；这一方面随着现代的线性群表示理论而变得更加清晰；由此人们便认识到，例如，魏尔斯特拉斯和克罗内克给出的行列式定义与由格拉斯曼的演算得出的定义实质上相同。

我们于是到达了现代，在那里，公理化方法和结构概念（起初被感受到，直到最近才被定义）使我们得以把直到那时还不可分解地混杂在一起的概念分离开来，把模糊的或潜意识的东西表述出来，以适当的普遍性去证明那些此前只在特殊情形中已知的定理。皮亚诺是公理化方法的创立者之一，也是最早按其真正价值赏识格拉斯曼著作的数学家之一，他早在 1888 年（{[}246 b{]}，第 IX 章）就给出了实数域上向量空间（不论是否有限维）的公理化定义，并以一种完全现代的记号给出了从这样一个空间到另一个空间的线性映射；稍后，平谢莱试图把这样构想出来的线性代数应用于函数论，其所取的方向诚然几乎没有显示出多产的迹象；至少他的观点使他得以在“拉格朗日伴随”中认出线性映射的转置的一个特殊情形：这一点不久便更清晰地出现，而且既出现在微分方程中，也出现在偏微分方程中，那是在希尔伯特及其学派关于希尔伯特空间及其分析应用的令人难忘的工作过程中。正是联系着后面这些研究，托普利茨 {[}309 a{]} 在（借助于坐标）引入实数域上最一般的向量空间的同时，作出了如下基本观察：行列式理论对于线性代数主要定理的证明是没有用处的，这就使得这些定理可以毫无困难地推广到无限维空间；他还指出，这样理解的线性代数当然适用于任意的交换基域。

另一方面，随着巴拿赫在 1922 年引入以他的名字命名的空间，\(^{3}\) 人们遇到了——诚然是在一个既属拓扑又属代数的问题中——与其对偶不同构的空间（见第~214~页以下）。在一个向量空间与其对偶空间之间，本来就不存在“典范”同构，也就是说由结构决定的同构，这一点长期以来就反映在同步与逆步的区分之中。尽管如此，似乎毫无疑问的是，一个空间与其对偶之间的区分，只是随着巴拿赫及其学派的工作才最终确立起来；也正是这一工作唤起了对余维数概念的兴趣。至于一个空间的向量子空间与其对偶空间的向量子空间之间的对偶性或“正交性”，它今天的表述方式显示出与伽罗瓦理论基本定理的现代表述、与局部紧阿贝尔群中所谓庞特里亚金对偶性的一种绝非偶然的类似；这可以追溯到韦伯，他在算术研究的过程中于 1866 年就有限群陈述了其基础 {[}327 d{]}；在伽罗瓦理论中，子群与子域之间的“对偶性”随戴德金和希尔伯特而成形；而向量子空间之间的正交性，则明显地一方面派生于射影几何中线性流形之间的对偶性，另一方面派生于欧氏空间或希尔伯特空间中完全正交流形的概念及其性质（它的名字即由此而来）。所有这些线索在当代汇合到代数学家如 E. 诺特、阿廷和哈塞、拓扑学家如庞特里亚金和惠特尼的手中，而且这两组人之间并非没有相互施加影响。

与此同时，人们着手进行一种批判性的审查，其意图是在每一点上消除那些并非真正不可缺少的假设，尤其是那些会妨碍某些应用的假设。由此便实现了在向量空间概念中以环代替域的可能性，而且，通过创立模的一般概念，可以同时处理这些空间、阿贝尔群、已被克罗内克、魏尔斯特拉斯、戴德金、施泰尼茨研究过的那些特殊的模，乃至带算子群，例如把若尔当—赫尔德定理应用于它们；与此同时，借助右模与左模的区分实现了向非交换性的过渡，而美国学派（韦德伯恩、迪克森）特别是德国学派（E. 诺特、阿廷）的代数理论的现代发展正导向这一区分。

最后，相当晚近才出现了我们必须在这里谈及的诸倾向中的最后一种：它以胚胎的形式包含于戴德金的定理（{[}79{]}，第 III 卷，第~29~页）——即交换域的任意多个自同构总是线性无关的——之中，伽罗瓦理论的线性化由阿廷 {[}7 a{]} 完成，随后很快被推广到交换域的任意扩张，甚至非交换域的扩张（{[}172 d{]}，第 VII 章）。

